\documentclass[aspectratio=169,11pt]{beamer}

\usetheme{Madrid}
\usecolortheme{seahorse}
\setbeamertemplate{navigation symbols}{}
\setbeamertemplate{footline}[frame number]

\usepackage[utf8]{inputenc}
\usepackage[T1]{fontenc}
\usepackage[spanish,es-nodecimaldot]{babel}
\usepackage{booktabs}
\usepackage{array}
\usepackage{amsmath}
\usepackage{tikz}
\usepackage{pgfplots}
\pgfplotsset{compat=1.18}

\definecolor{UdiGreen}{RGB}{31,92,78}
\definecolor{UdiBlue}{RGB}{30,68,118}
\definecolor{SoftGray}{RGB}{245,247,250}
\setbeamercolor{structure}{fg=UdiGreen}
\setbeamercolor{frametitle}{fg=white,bg=UdiGreen}
\setbeamercolor{title}{fg=white,bg=UdiGreen}

\title[SeLFIES en UDIS]{Adaptaci\'on de SeLFIES/RSB a UDIS}
\subtitle{Calculadora para estimar bonos de retiro y costo por edad}
\author{LAT4072 Seguridad Social y Pensiones}
\date{Datos de mercado: 12 de mayo de 2026}

\newcommand{\udis}{\textsc{UDIS}}
\newcommand{\mxn}{\textsc{MXN}}

\begin{document}

\begin{frame}
  \titlepage
\end{frame}

\begin{frame}{Objetivo ejecutivo}
  \begin{block}{Pregunta central}
    \centering
    \large ?`Cuantos bonos de retiro debe comprar una persona para asegurar una pension real objetivo?
  \end{block}
  \vspace{0.4cm}
  \begin{itemize}
    \item Adaptar el instrumento SeLFIES/RSB del articulo de Merton y Muralidhar a una unidad real mexicana: \udis.
    \item Usar UDIBONOS como proxy de la curva real de descuento.
    \item Calcular el numero de bonos y el costo actual para cualquier edad de entrada antes del retiro.
  \end{itemize}
\end{frame}

\begin{frame}{Idea del articulo}
  \begin{columns}[T]
    \begin{column}{0.53\textwidth}
      \begin{itemize}
        \item El problema no es solo acumular riqueza: es asegurar ingreso real durante el retiro.
        \item El articulo propone SeLFIES: \emph{Standard-of-Living indexed, Forward-starting, Income-only Securities}.
        \item El bono empieza a pagar al retiro, paga solo cupones y no requiere reinvertir pagos durante la etapa laboral.
      \end{itemize}
    \end{column}
    \begin{column}{0.42\textwidth}
      \begin{block}{Mensaje clave}
        El instrumento convierte una meta compleja de ahorro en una unidad simple:
        \[
          N = \frac{B}{c}
        \]
        donde $B$ es el ingreso anual objetivo y $c$ es el cupon por bono.
      \end{block}
    \end{column}
  \end{columns}
\end{frame}

\begin{frame}{Diseno del instrumento}
  \begin{center}
  \begin{tikzpicture}[x=0.95cm,y=0.8cm]
    \draw[->,thick] (0,0) -- (11.5,0) node[right]{tiempo};
    \draw[UdiBlue,very thick] (0,0.15) -- (5.2,0.15);
    \node[UdiBlue] at (2.6,0.55) {acumulacion};
    \draw[UdiGreen,very thick] (5.5,0.15) -- (10.7,0.15);
    \node[UdiGreen] at (8.2,0.55) {retiro};
    \draw[dashed] (5.4,-0.25) -- (5.4,2.8);
    \node[below] at (5.4,-0.25) {edad 65};
    \foreach \x in {5.6,6.1,...,10.2} {
      \draw[UdiGreen,thick] (\x,0) -- (\x,1.3);
    }
    \node[align=center] at (8,2.2) {cupones reales\\durante 20 anos};
    \draw[thick,red!70!black] (0.8,0) -- (0.8,-1.2);
    \node[align=center,red!70!black] at (1.9,-1.0) {precio\\hoy};
  \end{tikzpicture}
  \end{center}
  \begin{itemize}
    \item \textbf{Forward-starting:} no paga antes del retiro.
    \item \textbf{Income-only:} entrega cupones, sin principal final.
    \item \textbf{Real:} los pagos se expresan en una unidad que mantiene poder adquisitivo.
  \end{itemize}
\end{frame}

\begin{frame}{Adaptacion a Mexico}
  \begin{columns}[T]
    \begin{column}{0.48\textwidth}
      \begin{block}{Articulo}
        Indexacion ideal al consumo per capita para cubrir riesgo de estandar de vida.
      \end{block}
    \end{column}
    \begin{column}{0.48\textwidth}
      \begin{block}{Implementacion}
        Pagos denominados en \udis. Esto cubre inflacion, pero no necesariamente crecimiento del estandar de vida.
      \end{block}
    \end{column}
  \end{columns}
  \vspace{0.45cm}
  \begin{itemize}
    \item Se usa la UDI como unidad real de cuenta.
    \item Se usa la curva de UDIBONOS como aproximacion de tasas reales mexicanas.
    \item La app permite agregar crecimiento real adicional $g$ como sensibilidad.
  \end{itemize}
\end{frame}

\begin{frame}{Datos usados}
  \begin{columns}[T]
    \begin{column}{0.58\textwidth}
      \begin{table}
      \centering
      \begin{tabular}{rrrr}
        \toprule
        Plazo & Precio & Tasa & Tasa decimal \\
        \midrule
        3 anos  & 887.17 & 4.02\% & 0.0402 \\
        10 anos & 851.87 & 4.61\% & 0.0461 \\
        30 anos & 797.23 & 4.65\% & 0.0465 \\
        \bottomrule
      \end{tabular}
      \end{table}
    \end{column}
    \begin{column}{0.37\textwidth}
      \begin{block}{UDI}
        \centering
        Valor UDI usado:
        \[
          8.838416\ \mxn
        \]
      \end{block}
      \begin{itemize}
        \item Fecha de captura: 12-may-2026.
        \item Interpolacion lineal entre nodos.
        \item Extrapolacion plana despues de 30 anos.
      \end{itemize}
    \end{column}
  \end{columns}
\end{frame}

\begin{frame}{Supuestos del caso base}
  \begin{table}
  \centering
  \begin{tabular}{ll}
    \toprule
    Supuesto & Valor base \\
    \midrule
    Pension anual objetivo & 72,000 \udis \\
    Cupon anual por bono & 5 \udis \\
    Frecuencia de pagos modelada & Anual \\
    Edad de retiro & 65 anos \\
    Periodo de pago & 20 anos \\
    Primer pago & Al cumplir edad de retiro \\
    Capitalizacion & Anual \\
    Crecimiento real adicional $g$ & 0.0\% \\
    Rango de edades calculado & 16 a 55 anos, editable en la app \\
    \bottomrule
  \end{tabular}
  \end{table}
\end{frame}

\begin{frame}{Modelo de valuacion}
  Para una persona de edad $x$, con retiro a edad $R$, los pagos se modelan con frecuencia \textbf{anual}:
  \[
  \begin{aligned}
    n &= R - x
    &\qquad
    N &= \frac{B}{c} \\
    P_x &= \sum_{j=0}^{L-1} c(1+g)^j \cdot v(n+j)
    &\qquad
    v(t) &= \frac{1}{(1+r(t))^t} \\
    C_x &= N \cdot P_x
  \end{aligned}
  \]
  \begin{block}{Interpretacion}
    \begin{itemize}
      \item Cada flujo $c(1+g)^j$ representa el cupon de un ano de retiro.
      \item La edad no cambia el numero de bonos; cambia el precio hoy, porque los pagos quedan mas cerca o mas lejos.
    \end{itemize}
  \end{block}
\end{frame}

\begin{frame}{Resultado principal: bonos requeridos}
  \begin{center}
    \Huge 14,400 bonos
  \end{center}
  \vspace{0.2cm}
  \[
    \frac{72{,}000\ \udis}{5\ \udis\ \text{por bono}} = 14{,}400
  \]
  \vspace{0.3cm}
  \begin{itemize}
    \item Este numero es el mismo para cualquier edad.
    \item Lo que cambia con la edad es el costo de comprar hoy los flujos futuros.
    \item La app permite cambiar el cupon; si el cupon fuera 1 UDI, se requeririan 72,000 bonos.
  \end{itemize}
\end{frame}

\begin{frame}{Costo por edad}
  \begin{center}
  \begin{tikzpicture}
    \begin{axis}[
      width=0.86\textwidth,
      height=0.55\textheight,
      xlabel={Edad actual},
      ylabel={Costo total, millones de pesos},
      ymin=0,
      ymax=6,
      xmin=16,
      xmax=55,
      grid=major,
      axis line style={gray!70},
      tick label style={font=\small},
      label style={font=\small},
      line width=1pt,
      mark size=2pt
    ]
      \addplot[color=UdiGreen,mark=*] coordinates {
        (16,0.922) (20,1.106) (25,1.388) (30,1.742)
        (35,2.187) (40,2.746) (45,3.451) (50,4.336) (55,5.447)
      };
    \end{axis}
  \end{tikzpicture}
  \end{center}
  \vspace{-0.15cm}
  \begin{block}{Lectura}
    Empezar tarde encarece la misma pension objetivo: a 55 anos cuesta casi 5.9 veces lo que a 16 anos.
  \end{block}
\end{frame}

\begin{frame}{Resultados seleccionados}
  \begin{table}
  \centering
  \small
  \begin{tabular}{rrrrr}
    \toprule
    Edad & Anos al retiro & Precio/bono & Costo total & Costo total \\
    & & \udis & \udis & \mxn \\
    \midrule
    16 & 49 & 7.25 & 104,335 & 922,159 \\
    25 & 40 & 10.91 & 157,067 & 1,388,222 \\
    35 & 30 & 17.18 & 247,444 & 2,187,012 \\
    45 & 20 & 27.11 & 390,434 & 3,450,817 \\
    55 & 10 & 42.80 & 616,296 & 5,447,078 \\
    \bottomrule
  \end{tabular}
  \end{table}
  \begin{itemize}
    \item Todos los renglones compran la misma pension: 72,000 \udis{} al ano.
    \item El costo sube porque hay menos anos de descuento antes del primer pago.
  \end{itemize}
\end{frame}

\begin{frame}{Implementacion de la calculadora}
  \begin{columns}[T]
    \begin{column}{0.48\textwidth}
      \begin{block}{Datos}
        \begin{itemize}
          \item \texttt{curva\_original\_udibonos.csv}
          \item \texttt{curva\_real\_udis.csv}
          \item \texttt{valor\_udi.csv}
        \end{itemize}
      \end{block}
      \begin{block}{Script de curva}
        \texttt{calcular\_curva\_udibonos.py} convierte tasas porcentuales a decimales y genera la curva real usada por la app.
      \end{block}
    \end{column}
    \begin{column}{0.48\textwidth}
      \begin{block}{Motor}
        \texttt{selfies\_calculator.py}
        \begin{itemize}
          \item valida curva;
          \item interpola tasas;
          \item descuenta flujos;
          \item calcula bonos y costo por edad.
        \end{itemize}
      \end{block}
      \begin{block}{Interfaz}
        \texttt{app.py} permite cambiar edades, retiro, cupon, pension, tasa de crecimiento real y valor UDI.
      \end{block}
    \end{column}
  \end{columns}
\end{frame}

\begin{frame}{Consistencia con el articulo}
  \begin{table}
  \centering
  \small
  \begin{tabular}{p{0.32\textwidth}p{0.28\textwidth}p{0.30\textwidth}}
    \toprule
    Elemento & Articulo & Adaptacion en UDIS \\
    \midrule
    Inicio de pagos & Desde el retiro & Desde edad 65 \\
    Tipo de pago & Cupon real & Cupon en UDIS \\
    Principal final & No & No \\
    Plazo de retiro & Ejemplo de 20 anos & 20 anos base, editable \\
    Indexacion & Consumo per capita & UDI; $g$ opcional para sensibilidad \\
    Numero de bonos & Ingreso/cupon & 72,000/5 = 14,400 \\
    \bottomrule
  \end{tabular}
  \end{table}
\end{frame}

\begin{frame}{Limitaciones y extensiones}
  \begin{itemize}
    \item \textbf{Estandar de vida:} UDIS protegen inflacion, pero no el crecimiento del consumo per capita.
    \item \textbf{Longevidad:} la version base paga por plazo fijo; no es una renta vitalicia actuarial.
    \item \textbf{Curva:} solo hay nodos 3, 10 y 30 anos; se interpola y se extrapola plano para horizontes largos.
    \item \textbf{Mercado:} los precios observados son de UDIBONOS existentes; un SeLFIES real tendria precio de mercado propio.
  \end{itemize}
  \vspace{0.2cm}
  \begin{block}{Siguiente mejora natural}
    Incorporar una curva con mas nodos, sensibilidad de tasas y una opcion actuarial con probabilidades de supervivencia.
  \end{block}
\end{frame}

\begin{frame}{Conclusiones}
  \begin{enumerate}
    \item La adaptacion conserva la intuicion central de SeLFIES: medir retiro como ingreso real, no como monto acumulado.
    \item Con el caso base se requieren 14,400 bonos para comprar una pension de 72,000 \udis{} al ano.
    \item El costo presente depende fuertemente de la edad: comprar antes reduce mucho el costo actual.
    \item La implementacion es flexible: permite evaluar cualquier edad antes del retiro y ajustar supuestos clave.
  \end{enumerate}
  \vspace{0.35cm}
  \begin{center}
    \large La herramienta convierte la meta de pension en una cantidad concreta de bonos a comprar.
  \end{center}
\end{frame}

\end{document}
